Jacobson's commutativity theorem implies that a ring in which each element is potent is commutative, but the theorem does not by itself produce explicit equational derivations for fixed exponents. We study the first prime-power layer not covered by Brandenburg's explicit \(k=1\) and \(k=2\) arguments: characteristic-\(p\) rings satisfying \(x^{p^3}=x\). Brandenburg's period criterion proves the case \(\gcd(3,p^3-1)=1\), leaving exactly the primes \(p\equiv 1 \pmod 3\). For this remaining family we prove finite certificates for the two monomial constructive Wedderburn statements \(W_{p,3,T^p}\) and \(W_{p,3,T^{p^2}}\). The new step is a cyclic norm argument. On an irreducible cubic Frobenius component, a unit \(b\) satisfying \(ba=a^{p^i}b\) has \(b^3\) fixed; a finite norm section supplies \(u\) with \(u\sigma(u)\sigma^2(u)=b^3\). The normalized element \(u^{-1}b\) has cube \(1\), and when \(p\equiv 1 \pmod 3\), split-polynomial centrality in reduced rings forces it to be central, collapsing the Frobenius cycle. Exact symbolic checks verify the congruence split and the finite norm-component structure for representative primes. Thus, for every fixed prime \(p\), \(p^3\)-rings have equational commutativity certificates.
